\subsection{The parser/dispatcher boundary}
\label{sec:rc-boundary}

Every HTTP server in our sample is structured as a parser feeding
a dispatcher. The parser decodes wire bytes (chunked-TE framing,
header parsing, body length tracking). The dispatcher delivers
decoded body bytes into application code (ASGI handler, Go's
\fname{http.Handler}, Tower service, Node \fname{Readable},
Servlet \fname{getInputStream}, Erlang process mailbox, Pipelines
\fname{PipeReader}).

The per-chunk amplification class lives at the boundary between
those two abstractions. The parser is required by the wire format
to track chunk-level state; this is unavoidable and roughly
constant-cost per chunk in our sample (10s of nanoseconds to a
few microseconds in C, less than \SI{10}{\micro\second} in pure
managed code). The dispatcher's cost \emph{per delivery event}
varies wildly with the runtime: a Go goroutine wake costs a few
microseconds; an asyncio \fname{await receive} costs perhaps
\SI{5}{\micro\second}; a Node \fname{nextTick} drains the
microtask queue (superlinear in queue depth); a JVM thread
context switch costs $\sim$\SI{15}{\micro\second}.

Whether the dispatcher cost is paid \emph{per chunk} or
\emph{per recv batch} is the choice that determines per-chunk
amplification. None of the servers in our sample except Kestrel
pay it per recv-batch by default.

Figure~\ref{fig:arch-three-classes} sketches the three structural
classes as data-flow diagrams from kernel \texttt{recv()} to
application handler.

\begin{figure*}[t]
  \centering
  \begin{tikzpicture}[
    node distance=4mm and 6mm,
    every node/.style={font=\footnotesize},
    box/.style={draw, rounded corners=2pt, minimum height=6mm,
                minimum width=22mm, align=center, font=\footnotesize},
    parser/.style={box, fill=blue!12},
    dispatcher/.style={box, fill=orange!18},
    handler/.style={box, fill=green!18},
    cost/.style={font=\scriptsize\itshape, text=red!60!black, align=left},
    arr/.style={-{Latex[length=2mm,width=1.5mm]}, thick},
    chunkarr/.style={arr, draw=red!70!black},
    batcharr/.style={arr, draw=green!50!black, line width=0.9pt},
  ]

  % PANEL A: per-chunk (most servers)
  \begin{scope}[local bounding box=panelA]
    \node[parser]    (kA)  {kernel \texttt{recv()}};
    \node[parser,    below=of kA]   (pA1) {chunked decoder};
    \node[dispatcher,below=of pA1]  (dA)  {dispatcher\\(per-chunk)};
    \node[handler,   below=of dA]   (hA)  {application\\handler};
    \draw[arr] (kA) -- (pA1);
    \draw[chunkarr] (pA1) edge[bend left=45]   node[cost,right] {chunk 1} (dA);
    \draw[chunkarr] (pA1) edge[bend right=10]  node[cost,left]  {chunk 2} (dA);
    \draw[chunkarr] (pA1) edge[bend right=45]  node[cost,left]  {chunk N} (dA);
    \draw[chunkarr] (dA) -- node[cost,right] {N events: \\ N$\times$ task-switch} (hA);
    \node[above=2mm of kA,font=\small\bfseries] {(a) per-chunk};
  \end{scope}

  % PANEL B: batches-correctly (Kestrel)
  \begin{scope}[local bounding box=panelB, xshift=58mm]
    \node[parser]    (kB)  {kernel \texttt{recv()}};
    \node[parser,    below=of kB]   (pB1) {chunked decoder};
    \node[parser,    below=of pB1]  (pB2) {\fname{Pipe} ring buffer};
    \node[dispatcher,below=of pB2]  (dB)  {dispatcher\\(per-recv)};
    \node[handler,   below=of dB]   (hB)  {application\\handler};
    \draw[arr] (kB) -- (pB1);
    \draw[chunkarr] (pB1) edge[loop right, looseness=3.5, min distance=8mm] node[cost,right] {N$\times$ in pump} (pB1);
    \draw[arr] (pB1) -- (pB2);
    \draw[batcharr] (pB2) -- node[cost,right] {1 event: \\ 1$\times$ task-switch} (dB);
    \draw[arr] (dB) -- (hB);
    \node[above=2mm of kB,font=\small\bfseries] {(b) batches-correctly};
  \end{scope}

  % PANEL C: quadratic
  \begin{scope}[local bounding box=panelC, xshift=116mm]
    \node[parser]    (kC)  {kernel \texttt{recv()}};
    \node[parser,    below=of kC]   (pC1) {chunked decoder};
    \node[dispatcher,below=of pC1]  (dC)  {dispatcher\\(per-chunk)};
    \node[box, fill=red!18, below=of dC]   (qC)  {accumulator\\(grows with N)};
    \node[handler,   below=of qC]   (hC)  {application\\handler};
    \draw[arr] (kC) -- (pC1);
    \draw[chunkarr] (pC1) edge[bend left=45]  node[cost,right] {chunk 1} (dC);
    \draw[chunkarr] (pC1) edge[bend right=10] node[cost,left]  {chunk 2} (dC);
    \draw[chunkarr] (pC1) edge[bend right=45] node[cost,left]  {chunk N} (dC);
    \draw[chunkarr] (dC) -- node[cost,right] {N events $\times$ \\ O(M) walk} (qC);
    \draw[arr] (qC) -- (hC);
    \node[above=2mm of kC,font=\small\bfseries] {(c) quadratic};
  \end{scope}

  \end{tikzpicture}

  \caption{Three structural classes observed across the 27 HTTP/1.1 servers.
  (a)~\emph{per-chunk}: the chunked decoder dispatches one
  application receive event per chunk; the cost is N task-switches.
  (b)~\emph{batches-correctly} (Kestrel /
  \fname{System.IO.Pipelines}): the decoder loop pumps all
  chunks present in the recv buffer into a Pipe ring buffer and
  triggers one application receive event per recv batch; the
  per-chunk cost stays inside the pump and never crosses the
  \fname{async}/\fname{await} boundary.
  (c)~\emph{quadratic}: each per-chunk dispatch additionally
  walks an accumulator whose size grows with N, giving
  $\Theta(N^2)$ total work. In our sample, the three Node
  \fwid{http.Server}-based servers fit this $\Theta(N^2)$ pattern
  via the \fname{maybeReadMore}-driven microtask queue.}
  \label{fig:arch-three-classes}
\end{figure*}

\subsection{Three cost mechanisms}
\label{sec:rc-mechanisms}

We observe three distinct cost mechanisms across the 23
VULNERABLE-PER-CHUNK + 3 QUADRATIC servers.

\paragraph{Mechanism A: dispatcher-bound (most servers).}
The parser is fast; the cost is the per-chunk dispatch event.
Representative example: uvicorn + h11. Profile shows
$\sim$\SI{5}{\micro\second}/chunk, of which the h11 state-machine
step is $\sim$\SI{1}{\micro\second} and the
\fname{message\_event.set()} -> \fname{await receive()} round-trip
plus the ensuing async task switch is $\sim$\SI{4}{\micro\second}.

Most of our sample sits in this regime. Optimizing the parser
gives essentially no improvement; optimizing the dispatcher (by
batching) gives \emph{all} of the improvement.

\paragraph{Mechanism B: thread-handoff (Tomcat).}
On top of mechanism A, the thread-per-connection Servlet model
forces a Tomcat NIO Poller -> SynchronizedQueue -> worker thread
handoff per chunk. The
\fname{pthread\_cond\_signal} / \fname{ObjectMonitor::wait} traffic
visible in \fname{async-profiler} adds $\sim$\SI{14}{\micro\second}
on top of the underlying \fname{ChunkedInputFilter.parseChunkHeader}
cost. Total $\sim$\SI{21}{\micro\second}/chunk -- the worst
linear cell among the conventional servers at $N{=}\num{250000}$,
once the two high-constant outliers (the \fwid{bandit} iolist walk
at $\sim$\SI{138}{\micro\second} and the synthetic nginx-as-origin
cell at $\sim$\SI{114}{\micro\second}) are set aside.

\paragraph{Mechanism C: superlinear accumulator (Node).}
Some servers carry per-chunk \emph{state} that itself grows with
the number of chunks already received. When the per-chunk
operation has to inspect that state, the total work can become
quadratic. In our sample this manifests in the three Node
\fwid{http.Server}-based servers: the \fname{maybeReadMore} ->
\fname{process.nextTick} chain schedules a callback per chunk;
each scheduled callback drains some prefix of the microtask
queue; the queue grows with the number of pushed chunks, the
drain work per scheduled callback grows similarly, and total
work is $\Theta(N^2)$ (measured exponent $\sim 1.7$--$2.0$).

\fwid{Bandit} exhibits a structurally similar "per-chunk
operation walks accumulated state" pattern --
\fname{do\_read\_chunk!/4} calls \fname{erlang:iolist\_size/1}
on the accumulating iolist on every chunk (89.3\% of CPU in the
\fname{:eprof} sample) -- but in our $N$ range the cost does
\emph{not} grow quadratically: the measured exponent is
$\sim 1.05$ and per-chunk cost \emph{falls} with $N$
(155\,$\to$\,138\,\si{\micro\second}/chunk under Mode B),
indicating Bandit chunks the accumulator in a way that keeps the
per-chunk walk bounded over this range. We therefore treat
Bandit as a high-constant linear case, not a quadratic one.
(Adjacent prior art: CVE-2026-7790~\cite{cve_2026_7790_cowlib}
in \fwid{cowlib}'s chunk-size hex parser is a genuine quadratic
in the same library family.)

\subsection{Why Kestrel avoids both mechanisms}
\label{sec:rc-kestrel}

Kestrel's \fname{Http1ChunkedEncodingMessageBody.PumpAsync} drains
the entire socket-readable buffer into a \fname{System.IO.Pipelines}
\fname{Pipe} in a single synchronous loop. The loop transitions
through \fname{ParseChunkedPrefix} and \fname{ReadChunkedData}
states for every chunk in the buffer, writes the data bytes into
the \fname{Pipe} (chunk metadata is stripped), and only after the
buffer is fully drained does the pump call \fname{FlushAsync()}.
The application's \fname{PipeReader.ReadAsync} resumes
\emph{once} per pump, with a single \fname{ReadResult} containing
all decoded chunks' bytes coalesced.

This is mechanism-A elimination by construction. There is no
per-chunk dispatch event because the dispatcher (the
\fname{async/await} state machine of the handler) is only ticked
once per pump. There is no per-chunk accumulator walk because the
\fname{Pipe}'s ring buffer is amortized constant-cost-per-byte.
The result is a per-chunk cost (Section~\ref{sec:synth-classes})
dominated entirely by the synchronous parser loop, with no
runtime tax on top.

The Pipelines design is publicly documented and predates the
security framing of this bug class. Among HTTP/1.1 servers it is,
to our knowledge, the only production server in widespread use
that adopts this pattern as a default; four WebSocket
implementations (Section~\ref{sec:pf-ws}) independently arrive at
the same drain-before-dispatch design.

\subsection{Why batching is spec-compliant}
\label{sec:rc-spec}

We address a possible counterargument: that delivering one
application event per chunk is required by the relevant
specifications.

\begin{itemize}
  \item RFC 9112~\cite{rfc9112} defines the wire format. It
    does not specify any application-API granularity.
  \item The ASGI specification~\cite{asgi_spec} says a server
    "may send a single \fname{http.request} event with the entire
    body" or "may send the body as multiple
    \fname{http.request} events". The spec explicitly permits
    coalescing.
  \item PEP 3333 (WSGI)~\cite{wsgi_pep3333} defines
    \fname{environ['wsgi.input']} as a file-like object read by
    the application; the server is free to back it with any
    buffer it likes (Tempfile in Puma; \fname{BytesIO} in
    waitress).
  \item Java's Servlet API exposes \fname{getInputStream()};
    chunking is invisible to the handler.
  \item Node's \fname{IncomingMessage} extends
    \fname{Readable}, which is documented as preserving chunk
    boundaries when called in \fname{Buffer} mode, but the docs
    do not prohibit coalescing inside \fname{push()}.
\end{itemize}

In every case the spec permits the batching that Kestrel does.
The choice not to batch is implementation-internal, not
spec-mandated.
